\documentclass[a4paper,fleqn]{cas-sc}
\usepackage[numbers,sort&compress]{natbib}

\shorttitle{Short title of the article}
\shortauthors{F. Author et al.}

\begin{document}

\title[mode=title]{Title of the Study: A Prospective Cohort of Adults in Primary Care}

\author[1]{First Author}[orcid=0000-0000-0000-0000]
\cormark[1]
\ead{first.author@example.org}
\credit{Conceptualization, Methodology, Writing -- original draft}

\author[1,2]{Second Author}
\credit{Formal analysis, Writing -- review \& editing}

\author[2]{Third Author}
\credit{Data curation, Investigation}

\affiliation[1]{organization={Department of Medicine, University Hospital},
  city={City},
  postcode={00000},
  country={Country}}
\affiliation[2]{organization={School of Public Health, University},
  city={City},
  country={Country}}

\cortext[1]{Corresponding author.}

\begin{abstract}
\noindent\textbf{Background:} State why the question matters and what is unknown.\\
\textbf{Methods:} Describe the design, setting, participants and main outcome.\\
\textbf{Results:} Give the main findings with effect sizes and 95\% confidence intervals.\\
\textbf{Conclusions:} Say what the findings mean for practice or research.
\end{abstract}

\begin{highlights}
\item First key finding in one sentence.
\item Second key finding in one sentence.
\item What the findings mean for patients or clinicians.
\end{highlights}

\begin{keywords}
cohort study \sep primary care \sep risk factors \sep outcomes
\end{keywords}

\maketitle

\section{Introduction}
Summarise what is known and the gap this study addresses. Observational
studies are reported following STROBE~\cite{vonelm2007strobe}, and the study
was conducted in line with the Declaration of Helsinki~\cite{wma2013helsinki}.

\section{Methods}
\subsection{Study design and participants}
Describe the setting, dates, eligibility criteria and how participants were
recruited.

\subsection{Outcomes}
Define the primary and secondary outcomes and how they were measured.

\subsection{Statistical analysis}
Describe the models used. For example, the hazard of the outcome was modelled as
\begin{equation}
  h(t \mid x) = h_0(t) \exp\bigl(\beta^\top x\bigr),
\end{equation}
with $x$ the exposure and adjustment variables.

\section{Results}
Table~\ref{tab:baseline} shows the participants' characteristics at baseline,
and Figure~\ref{fig:flow} how they were selected.

\begin{table}[pos=htbp]
  \centering
  \caption{Baseline characteristics of the participants.}
  \label{tab:baseline}
  \begin{tabular}{lrr}
    \toprule
    Characteristic & Exposed ($n=412$) & Unexposed ($n=398$) \\
    \midrule
    Age, years, mean (SD) & 58.4 (11.2) & 57.9 (10.8) \\
    Women, $n$ (\%) & 221 (53.6) & 207 (52.0) \\
    Diabetes, $n$ (\%) & 64 (15.5) & 49 (12.3) \\
    \bottomrule
  \end{tabular}
\end{table}

\begin{figure}[pos=htbp]
  \centering
  \includegraphics[width=0.5\linewidth]{example-image}
  \caption{Selection of participants. Upload an image and use its file name here.}
  \label{fig:flow}
\end{figure}

\section{Discussion}
Interpret the findings, compare them with earlier studies, and discuss
strengths and limitations.

\section{Conclusions}
State the main conclusion in a few sentences.

\section*{Ethics approval}
Name the ethics committee and the approval number, and how consent was obtained.

\section*{Funding}
List the funders and grant numbers, or state that there was no specific funding.

\section*{Declaration of competing interest}
The authors declare no competing interests.

\section*{Data availability}
Say where the data can be obtained, and under which conditions.

\printcredits

\bibliographystyle{cas-model2-names}
\bibliography{references}

\end{document}
